% Part: computability
% Chapter: computability-theory
% Section: universal-part-function

\documentclass[../../../include/open-logic-section]{subfiles}

\begin{document}

\olfileid[pt]{cmp}{thy}{uni}
\olsection{A função parcialmente computável universal}

\begin{thm}
\ollabel{thm:univ-comp}
Existe uma função parcialmente computável universal~$\fn{Un}(e,x)$. Em
outras palavras, existe uma função $\fn{Un}(e,x)$ tal que:
\begin{enumerate}
\item $\fn{Un}(e,x)$ é parcialmente computável.
\item Se $f(x)$ é uma função parcialmente computável qualquer, então
existe um número natural $e$ tal que
$f(x) \simeq \fn{Un}(e,x)$ para todo~$x$.
\end{enumerate}
\end{thm}

\begin{proof}
Seja $\fn{Un}(e,x) \simeq U(\umin{s}{T(e,x,s)})$, em que $U$ e $T$ são
como no teorema da forma normal de Kleene
(\olref[nfm]{thm:normal-form}).
\end{proof}

\begin{explain}
Essa é apenas uma maneira precisa de dizer que temos uma enumeração
efetiva das funções parcialmente computáveis; a ideia é que, se
escrevermos $f_e$ para a função definida por
$f_e(x) = \fn{Un}(e,x)$, então a sequência $f_0$, $f_1$, $f_2$, \dots
inclui todas as funções parcialmente computáveis, com a propriedade de
que $f_e(x)$ pode ser calculada ``uniformemente'' em $e$ e~$x$. Para
simplificar, estamos usando uma função binária que é universal para
funções unárias, mas, codificando sequências de números, podemos
generalizar isso facilmente para mais argumentos. Por exemplo, observe
que, se $f(x,y,z)$ é uma função recursiva parcial de $3$ argumentos,
então a função $g(x) \simeq f((x)_0, (x)_1, (x)_2)$ é uma função
recursiva unária.
\end{explain}

\end{document}

